\section{実験}\label{sec:experiments}

\subsection{実験設定}

合成データを用いてDP-DIMフレームワークを評価する。
データ生成では、$n$トランザクション（$n = 300$、スケーラビリティ実験では100--2000）、$|\mathcal{I}| = 20$アイテムを使用し、指定区間$[s_d, e_d]$にアイテムセット$X_d$を出現確率0.85で注入する（背景出現確率0.05）。
ウィンドウサイズ$W = 20$、閾値$\theta = 5$、最大アイテムセットサイズ$K = 3$とする。
評価指標としてヤッカード類似度（Jaccard）、F1スコア（位置ベースのprecision/recall）を用いる。
各条件で10シード（42--51）の平均を報告する。
非DPの密集区間検出結果をオラクル基準（ground truth）として使用する。

\subsection{E1: プライバシー・精度トレードオフ}

$\varepsilon$を$\{0.1, 0.5, 1.0, 2.0, 5.0, 10.0\}$で変化させ、非DPの密集区間を真値としてJaccardおよびF1を測定した（表\ref{tab:e1}）。

\begin{table}[t]
\centering
\caption{E1: $\varepsilon$に対するJaccardおよびF1（$W=20$, $\theta=5$, 10シード平均）}
\label{tab:e1}
\begin{tabular}{crrrr}
\toprule
$\varepsilon$ & Jaccard (平均) & Jaccard (標準偏差) & F1 (平均) & F1 (標準偏差) \\
\midrule
0.1  & 0.429 & 0.031 & 0.598 & 0.030 \\
0.5  & 0.430 & 0.031 & 0.599 & 0.029 \\
1.0  & 0.431 & 0.030 & 0.600 & 0.029 \\
2.0  & 0.433 & 0.029 & 0.602 & 0.028 \\
5.0  & 0.439 & 0.030 & 0.608 & 0.029 \\
10.0 & 0.449 & 0.029 & 0.617 & 0.027 \\
\bottomrule
\end{tabular}
\end{table}

$\varepsilon$の増大に伴いJaccard・F1ともに単調に向上し、$\varepsilon = 10.0$で最大Jaccard 0.449、F1 0.617を達成した。
$\varepsilon = 1.0$でもF1 0.60を維持しており、実用的なプライバシーレベルで合理的な精度が得られることを示している。

\textbf{精度低下率の分析}：非DP手法のJaccardは0.461（同一合成データ、ノイズなし）であり、$\varepsilon = 1.0$でのDP手法の精度低下率は$(0.461 - 0.431) / 0.461 = 6.5\%$に留まる。$\varepsilon = 0.1$でも低下率は6.9\%であり、ウィンドウカウントの感度が1と小さいことがノイズ耐性の根源である。この結果は、密集区間マイニングがDPと相性の良い問題構造を持つことを示唆している。

\subsection{E2: メカニズム比較}

ラプラスメカニズム、ガウスメカニズム（$\delta = 10^{-5}$）、SVT（$c_\text{max} = 50$）の3手法を$\varepsilon = 2.0$で比較した（表\ref{tab:e2}）。

\begin{table}[t]
\centering
\caption{E2: メカニズム比較（$\varepsilon = 2.0$, 10シード平均）}
\label{tab:e2}
\begin{tabular}{lrrrr}
\toprule
メカニズム & Jaccard (平均) & Jaccard (標準偏差) & F1 (平均) & 実行時間 (ms) \\
\midrule
Laplace  & 0.348 & 0.083 & 0.516 & 0.31 \\
Gaussian ($\delta\!=\!10^{-5}$) & 0.311 & 0.049 & 0.475 & 0.37 \\
SVT ($c_\text{max}\!=\!50$) & 0.233 & 0.089 & 0.378 & 0.11 \\
非DP（オラクル） & 0.461 & 0.025 & 0.631 & 0.28 \\
\bottomrule
\end{tabular}
\end{table}

ラプラスメカニズムが最も高いJaccard（0.348）を達成した。
ガウスメカニズムはやや低いが標準偏差が小さく安定している。これは$(\varepsilon, \delta)$-DPの緩和が一定の精度改善を与えつつ、ガウスノイズの軽い裾が外れ値を抑制するためである。
SVTは最も高速であるが、ノイズ閾値との比較という性質上、区間境界の精度がやや低い。ただし、クエリ数$|L|$に依存しない予算消費は大規模データへのスケーリングにおいて有利である。
非DPオラクルとの比較により、$\varepsilon = 2.0$でラプラスメカニズムの精度低下は24.5\%であることがわかる。

\subsection{E3: スケーラビリティ}

トランザクション数$n$を$\{100, 200, 500, 1000, 2000\}$で変化させ、実行時間を測定した（表\ref{tab:e3}）。

\begin{table}[t]
\centering
\caption{E3: スケーラビリティ（$\varepsilon = 1.0$, 3シード平均）}
\label{tab:e3}
\begin{tabular}{rrr}
\toprule
$n$ & 実行時間 (ms, 平均) & 実行時間 (ms, 標準偏差) \\
\midrule
100   & 2.2  & 0.1 \\
200   & 4.5  & 0.0 \\
500   & 12.2 & 0.2 \\
1000  & 24.9 & 0.9 \\
2000  & 51.2 & 1.1 \\
\bottomrule
\end{tabular}
\end{table}

実行時間は$n$に対してほぼ線形にスケールしており（回帰係数$R^2 = 0.999$）、$n = 2000$でも51 msと実用的な速度である。
DPのオーバーヘッド（ノイズ生成）は全体の実行時間に対して無視できる水準であり、計算コストのボトルネックはウィンドウカウント計算にある。

\subsection{E4: 閾値安定性}

$\varepsilon$と信頼度$p$に対する安定性マージン$\gamma$を算出した（表\ref{tab:e4}）。これは定理\ref{thm:stability}に基づき、真のカウントが閾値から$\gamma$以上離れていれば、確率$p$以上で正しい密集判定が保証されることを意味する。

\begin{table}[t]
\centering
\caption{E4: 安定性マージン$\gamma$（$\Delta_W = 1$, 単一クエリ）}
\label{tab:e4}
\begin{tabular}{crrrr}
\toprule
$\varepsilon$ & $p = 0.80$ & $p = 0.90$ & $p = 0.95$ & $p = 0.99$ \\
\midrule
0.5 & 3.22 & 4.61 & 5.99 & 9.21 \\
1.0 & 1.61 & 2.30 & 3.00 & 4.61 \\
2.0 & 0.80 & 1.15 & 1.50 & 2.30 \\
5.0 & 0.32 & 0.46 & 0.60 & 0.92 \\
\bottomrule
\end{tabular}
\end{table}

$\varepsilon = 2.0$、信頼度$p = 0.95$で安定性マージンは1.50であり、真のカウントが閾値から2以上離れていれば95\%の信頼度で正しい判定が保証される。
$\varepsilon = 5.0$では$\gamma = 0.60$（$p = 0.95$）となり、閾値近傍のカウントでもほぼ安定した判定が可能である。
実用的には$\theta \geq 5$の設定ではカウント値が0または$\theta$以上に二極化するケースが多く、$|c - \theta| \geq 2$が大半のウィンドウ位置で成立するため、$\varepsilon \geq 1.0$で十分な安定性が得られる。

\subsection{E5: 予算合成}

逐次合成と高度合成定理による総$\varepsilon$消費を比較した（表\ref{tab:e5}、$\varepsilon_0 = 0.01$）。

\begin{table}[t]
\centering
\caption{E5: 逐次合成 vs 高度合成（$\varepsilon_0 = 0.01$, $\delta' = 10^{-5}$）}
\label{tab:e5}
\begin{tabular}{rrrr}
\toprule
クエリ数$k$ & 逐次合成$\varepsilon$ & 高度合成$\varepsilon$ & 削減率 \\
\midrule
10  & 0.10 & 0.15 & $-52.7$\% \\
50  & 0.50 & 0.34 & 31.1\% \\
100 & 1.00 & 0.49 & 51.0\% \\
500 & 5.00 & 1.12 & 77.5\% \\
\bottomrule
\end{tabular}
\end{table}

クエリ数が少ない場合（$k = 10$）は高度合成の$\sqrt{k\ln(1/\delta')}$項のオーバーヘッドが支配的で逐次合成が有利であるが、$k \geq 50$では高度合成が優位となる。
$k = 500$で77.5\%の予算削減を達成しており、多数の候補アイテムセットを探索する実用シナリオで実質的な利点がある。

\subsection{E6: パラメータ感度分析}

ウィンドウサイズ$W$と閾値$\theta$がDP精度に与える影響を分析した（表\ref{tab:e6}）。$\varepsilon = 1.0$で固定し、$W \in \{10, 20, 40\}$、$\theta \in \{3, 5, 10\}$を変化させた。

\begin{table}[t]
\centering
\caption{E6: パラメータ感度分析（$\varepsilon = 1.0$, $n = 300$, 10シード平均F1）}
\label{tab:e6}
\begin{tabular}{crrrr}
\toprule
& & \multicolumn{3}{c}{$\theta$} \\
\cmidrule(lr){3-5}
$W$ & 非DP F1 & $\theta = 3$ & $\theta = 5$ & $\theta = 10$ \\
\midrule
10  & 0.583 & 0.541 & 0.572 & 0.580 \\
20  & 0.631 & 0.558 & 0.600 & 0.624 \\
40  & 0.672 & 0.587 & 0.639 & 0.667 \\
\bottomrule
\end{tabular}
\end{table}

$\theta$が大きいほどDP精度が非DPに接近する傾向が観察される。これは閾値安定性マージンに関連しており、$\theta$が大きいとき、密集ウィンドウのカウント値が閾値から十分に離れるため、ノイズによる誤判定が減少する。
$W$が大きいほど精度が向上するのは、ウィンドウ内カウントの絶対値が増加し、信号対ノイズ比（SNR）が改善するためである。
これらの結果から、$\theta / \gamma \geq 2$（すなわち$\theta \geq 2\gamma$）を設計指針とすることが望ましい。
